\chapter{Introducción}
\label{ch:introduccion}
\section{Motivación y planteamiento del problema}
Un aerogenerador convierte parte de la energía del viento en energía mecánica mediante un rotor conectado al sistema de generación eléctrica. Desde el punto de vista estructural, esta máquina reúne una torre esbelta, palas flexibles y componentes cuya deformación puede ser pequeña frente al movimiento del conjunto. La rotación del rotor introduce una diferencia fundamental respecto de una estructura estacionaria: las partículas de las palas poseen velocidad y aceleración incluso cuando no existe vibración relativa.

La respuesta frente a una excitación externa depende tanto de las propiedades de masa y rigidez como del estado de operación. Una aceleración del terreno se transmite desde la base hacia la torre y el rotor; una carga aerodinámica actúa sobre una superficie cuya posición y orientación cambian durante el movimiento. Por esta razón, el análisis de las vibraciones exige describir conjuntamente las deformaciones estructurales, el giro del rotor y la transmisión de fuerzas y momentos entre componentes.

El problema abordado consiste en formular las ecuaciones de la dinámica estructural de un aerogenerador de eje horizontal sometido a cargas aerodinámicas y movimiento de base, conservando los efectos de su rotación prescrita. La dificultad no se limita a construir matrices de masa y rigidez. Es necesario establecer qué movimientos son conocidos, cuáles constituyen incógnitas, en qué sistema de referencia se expresan las variables y alrededor de qué estado se estudian las perturbaciones. Estas decisiones determinan el significado físico de las frecuencias naturales, los modos de vibración y las fuerzas equivalentes que intervienen en el modelo.

\section{Antecedentes y fundamentos del modelo}
La descripción de un cuerpo deformable parte de una configuración de referencia y de una aplicación que asigna la posición actual de cada partícula. Este enfoque permite distinguir desplazamiento, rotación y deformación y constituye la base de la mecánica del continuo desarrollada por Holzapfel \cite{holzapfel}. Para un sistema compuesto por varias piezas, la dinámica de sistemas multicuerpo proporciona las relaciones de compatibilidad y las transformaciones entre sistemas de referencia \cite{shabana}. Ambos enfoques son complementarios: el primero describe la deformación dentro de un componente y el segundo organiza el movimiento relativo entre componentes.

En el estudio de rotores eólicos, Gebhardt y colaboradores \cite{gebhardt2008} emplean sistemas de referencia asociados al cubo y a las palas para formular su cinemática y relacionarla con las cargas aerodinámicas. La identificación explícita de estos sistemas resulta especialmente útil cuando se distinguen el giro del rotor, el ángulo de paso y la orientación de cada pala. En esta tesis se conserva esa organización geométrica, mientras que el movimiento nominal se restringe a una rotación de velocidad constante alrededor del eje del rotor.

La representación estructural puede establecerse con diferentes niveles de detalle. Las formulaciones de vigas geométricamente exactas permiten tratar rotaciones finitas y deformaciones de componentes esbeltos \cite{simo1985}. BeamDyn aplica este enfoque al análisis de palas de aerogeneradores \cite{wang2017}. La comparación de elementos de viga y de lámina realizada por Faccio Júnior y colaboradores \cite{faccio2019} sitúa ambas representaciones en el contexto de la modelación de palas. La elección de una teoría de viga orienta el análisis hacia desplazamientos globales y resultantes internas; la descripción detallada de las paredes y uniones requiere variables adicionales.

Los trabajos de Maza y colaboradores \cite{maza2017} y Turello y colaboradores \cite{turello2023} proporcionan antecedentes sobre los fenómenos aeroelásticos y la respuesta de torres eólicas frente a acciones sísmicas. La formulación aquí desarrollada se concentra en la estructura y en la incorporación mecánicamente consistente de sus cargas. La teoría de vibraciones de Meirovitch \cite{meirovitch} permite organizar las contribuciones del movimiento prescrito y obtener las ecuaciones de pequeñas perturbaciones. La discretización espacial se apoya en el método de elementos finitos \cite{reddy}.

\section{Objetivo de la investigación}
El objetivo general es desarrollar un modelo de la dinámica estructural de un aerogenerador de eje horizontal que incorpore los efectos giroscópicos de la rotación del rotor y las cargas aerodinámicas dependientes de la configuración, junto con la excitación sísmica de la base. El modelo está dirigido al estudio de las vibraciones, la estabilidad y la respuesta dinámica en torno a un estado de operación.

Este objetivo requiere construir una descripción cinemática común para los cuerpos rígidos y flexibles; obtener sus propiedades de masa y rigidez con una convención uniforme de ejes y puntos de referencia; formular la discretización y las conexiones entre componentes; y deducir las ecuaciones de movimiento y su linealización. La comparación de representaciones estructurales exige, además, distinguir los cambios debidos a la aproximación espacial de aquellos asociados a hipótesis físicas diferentes.

\section{Organización del desarrollo}
La exposición parte del sistema mecánico y sus hipótesis. La cinemática establece las configuraciones, los sistemas de referencia y el movimiento; después se definen las propiedades geométricas, de masa y de rigidez, y la discretización mediante elementos finitos. Las ecuaciones dinámicas incorporan el movimiento prescrito y las cargas para describir el equilibrio de operación y sus perturbaciones. Cada magnitud se define antes de discretizarla.

